\subsection{Distributional evaluation and loss scaling}
\label{sec:distribution-followups}
The strict engineering gate compares each generated future with one realized
reference. D024 and D030 show that even simulator samples with valid hidden
physics can fail it. D031 therefore adds a descriptive multi-reference energy
score on pose displacements in native pixels. For $K$ sampled trajectories and
one reference $y$, the fair finite-ensemble form is
\[
 \widehat{ES}(X,y)=\frac{1}{K}\sum_i d(x_i,y)
 -\frac{1}{2K(K-1)}\sum_{i\ne j}d(x_i,x_j),
\]
averaged over the available references of each condition. Here $d$ is the
Euclidean trajectory distance divided by the square root of the number of
frames. Under independent sampling from the forecast distribution, the fair
form removes the finite-ensemble bias; the underlying energy score is proper
\cite{gneiting,ferro}. Four samples and at most four references still give noisy
per-condition estimates. Comparisons use paired scene-group intervals, and the
same four sampled worlds are shared across conditions. This does not establish
calibrated uncertainty or change the original gate.

D031 reproduces the recorded D028 scorer fields, with same-world ADE agreeing
within $1.5\times10^{-10}$. On that corpus D020 scores 1.105 pixels against
1.060 for the prior-matched simulator; the paired difference interval is
$[0.019,0.074]$. The 0.15-pixel hull baseline is within sampling uncertainty of
the simulator on validation, D028 and D033. Every evaluated learned checkpoint
remains worse than the simulator on this distributional measure. Matching a
single-reference ADE is therefore insufficient evidence of a matched forecast
distribution.

D032 investigates why the optional geometric and contact-dynamics terms had
little effect. In matched retraining, the D020 recipe's weighted geometry terms
contributed about 0.5\% of total training loss, while the D029 contact-dynamics
terms contributed about 0.02\%. The screen scales the D020 geometry terms by
16, 128 and 1024, and the D029 contact-dynamics terms by 128 and 1024. The
original-validation rule selects at most one scale per family: rest-anchored
$\times16$ and direct contact-dynamics $\times1024$. Increasing the rest-family
weights reduces penetration but worsens same-world ADE and energy score at
every tested scale. Retrained controls match the original model and EMA weights
bitwise. The initial scoring attempt stopped at an invalid assumption that
zero-force images must be fixed for direct models. The rerun records their
actual drift without changing any rule, margin or arm.

D033 evaluates the selections and six E10 checkpoints on ten new ID scene
groups. The projection passes its declared checks for all ten checkpoints.
The direct contact-dynamics selection improves at training seed 19, but D034
fails to replicate that complete rule at either seed 20 or 21: wall overshoot
and lost tracks remain. Penetration alone reaches zero at all three seeds.
The raw losses do not establish a general geometry improvement.

D035 tests the scaled direct losses together with projection on another ten
new ID groups, generated after that comparison was registered. All three
training seeds pass the combined rule. Same-world selected ADE falls by
8 to 10\%, or 0.03 to 0.05 native pixels, with zero projected geometry violations
or lost tracks. The comparison was chosen after seeing D034; the failed raw
replication is retained. Selected first-reference ADE and energy score do not
improve at every seed, and every learned run remains worse than the simulator
on energy score. This supports the limited hybrid comparison, not an overall
learned-physics or OOD claim.

\subsection{Candidate selection: a negative accuracy result}
\label{sec:candidate-selection}
D036 generates 16 candidates from one frozen E11 checkpoint on the already
opened D018 ID corpus, then selects four using visible hull violations,
stationarity, force opposition and trajectory diversity. The selector sees no
reference future or true hidden physics. Its fair control is the exact
expectation over uniform four-of-16 subsets, rather than a favorable or
unfavorable choice of the first four seeds.

\begin{table}[!ht]
\centering
\begin{tabular}{lrrr}
\hline
D036, 240 returned clips & Uniform expectation & Selector & Oracle \\
\hline
Selected ADE & 0.02621 & 0.02760 & 0.01245 \\
Complete clips & 237.25 & 240 & 240 \\
Penetrating clips & 19.0 & 0 & 16 \\
Boundary-clipped clips & 27.75 & 14 & 29 \\
\hline
\end{tabular}
\caption{D036 candidate selection. Uniform ADE averages complete candidates;
its missing-clip expectation is reported separately. The oracle uses the
reference future and is unavailable at inference. Geometry counts use the
same hull definitions as the selector and are not independent evidence.}
\label{tab:selection-negative}
\end{table}

Selection raises completeness but worsens ADE by 5.28\% against the fair
control. The oracle identifies substantial candidate-pool headroom that this
hand-written score cannot recover. Sixteen candidates also require four times
the model sampling of four direct outputs. The selector is not promoted to
the public sampler. The accepted D028 projection is exposed separately as an
explicit optional correction, preserving the ordinary seeded sampling path.

\subsection{Dynamics residual and momentum correction}
\label{sec:dynamics-residual}
D038 adds two separate optional components. A training residual matches
finite-difference free-motion acceleration to the known force divided by
sampled mass, excluding object and wall contacts where impulses make that
target incomplete. An inference block corrects post-force linear momentum
using the known applied impulse and sampled masses. Before predicted contact,
the correction acts on the selected object; after contact, it adjusts the
centre-of-mass motion while preserving relative velocity. Rotation is
unchanged. Trajectories approaching a wall are skipped because wall impulse
is unobserved. The inference order is neural sampling, momentum correction,
then the accepted hull projection. These are explicitly labelled hybrid
corrections, rather than evidence that the model learned the enforced laws.

The matched seed-20 pair uses the same E7 teacher, initialization, private RNG
and 8,000 updates; only the dynamics-residual weight differs, zero versus four.
The fresh seed-980000 corpus contains ten ID scene groups, 60 conditions and
240 generated clips per arm. All 355 stored simulator clips replay exactly.
The rule was committed locally before generation, but the registration and
results first reached GitHub together. Remote history therefore does not
independently establish the registration ordering.

\begin{table}[!ht]
\centering
\begin{tabular}{lrr}
\hline
D038 fresh ID & Projected control & Full hybrid \\
\hline
Complete clips & 240/240 & 240/240 \\
Penetrating / boundary clips & 0 / 0 & 0 / 0 \\
Zero-force / passive drift & 0 / 0 & 0 / 0 \\
Same-world selected ADE & 0.00699 & 0.00582 \\
First-reference pose ADE & 0.02771 & 0.02665 \\
Free-motion dynamics residual & 0.00564 & 0.00101 \\
Energy score, pixels & 1.2605 & 1.2539 \\
RGB contact F1 & 0.228 & 0.213 \\
\hline
\end{tabular}
\caption{D038 full hybrid: dynamics-residual training, momentum correction and
projection. ADE uses normalized coordinates. First-reference pose ADE is
distinct from RGB-tracked ADE. Zero geometry violations are measured on
observed hulls and do not establish complete physical validity.}
\label{tab:dynamics-residual}
\end{table}

All seven comparison checks pass. Same-world ADE falls by 16.8\%, with a
paired scene interval of $[-0.00149,-0.00084]$. The first-reference pose ADE
gain is 3.8\%. The recorded momentum-conservation check is separate: the block
enforces it by construction. Same-world ADE also uses the same sampled
masses supplied to that block, making the component comparison essential.

\begin{table}[!ht]
\centering
\begin{tabular}{lr}
\hline
D038 component & Same-world selected ADE \\
\hline
Control plus projection & 0.00699 \\
Add momentum correction & 0.00618 \\
Residual training plus projection, no momentum block & 0.00688 \\
Residual plus momentum plus projection & 0.00582 \\
\hline
\end{tabular}
\caption{D038 component attribution. The explicit momentum block accounts for
about two thirds of the aggregate gain and all of the wall-family gain.}
\label{tab:dynamics-components}
\end{table}

The full pipeline's collision-family ADE rises by 1.17\%, and RGB contact F1
falls. The momentum block alone worsens collision motion; the residual
recovers most of that loss. Without the block, residual training reduces its
trained dynamics error by 80.7\% and raises contact F1 from 0.228 to 0.278,
but gives only a small same-world ADE improvement. One confirmatory training
seed and ten ID groups support no collision-specific or OOD claim. TC2 A40
job 53377 independently retrains the pair and reproduces the principal values
to about four significant figures on a regenerated corpus; this is a
cross-platform reproduction, not another training-seed replication.

D039 retains a separate exploratory collision-impulse screen. Both seed-22
arms train for 2,000 updates with the dynamics residual and use momentum
correction plus projection; only the collision-residual weight differs, zero
versus 1024. On the already opened D038 corpus, collision same-world ADE
changes from 0.00683 to 0.00606, with interval $[-0.00117,-0.00043]$, while
contact F1 changes from 0.168 to 0.238. Both retain 240 complete clips without
measured geometry violations or drift. First-reference pose ADE becomes
slightly worse, and raw neural collision ADE worsens by 3.2\%. A shortened,
single-seed screen on an opened corpus is not confirmation or evidence of a
standalone learned-collision improvement. The default-zero loss remains
experimental. D037 was unrecorded development work and contributes no numeric
evidence to this report.
